\SourceLeaf{309}%
werth, das Verfahren anzugeben, nach welchem diese Entwicklung beliebig weit getrieben werden kann.

Ich erlaube mir an den dort gebrauchten Bezeichnungen einige Abänderungen, theils des bequemern Drucks wegen, theils um den verschiedenen Bezeichnungen in den einzelnen Theilen der gegenwärtigen Abhandlung etwas mehr Symmetrie geben zu können. Zunächst bedeute hier
\begin{list}{}{\setlength{\leftmargin}{1.4em}\setlength{\itemindent}{-1.4em}\setlength{\itemsep}{0pt}\setlength{\parsep}{0pt}}
\item \(r\) die Entfernung der beiden Punkte von einander, den Halbmesser der Kugel als Einheit angenommen.
\item \(B+\frac12b\) und \(B-\frac12b\) die Breite am ersten und zweiten Endpunkte von \(r\).
\item \(T+\frac12t\) und \(T-\frac12t\pm180^0\) das Azimuth des zweiten und ersten Endpunkts resp. vom ersten und zweiten aus.
\item \(l\) den Längenunterschied.
\end{list}
\noindent Es wird angenommen, dass das Azimuth von Süden nach Westen zu gezählt und \(l\) als positiv betrachtet wird, wenn der zweite Punkt westlicher liegt als der erste.

Es soll ferner gesetzt werden
\[
\begin{aligned}
\beta&=r\cos T\\
\tau&=r\sin T\mathord{.}\operatorname{tang}B\\
\lambda&=\frac{r\sin T}{\cos B}
\end{aligned}
\]
welche Grössen dasselbe ausdrücken, was im 17. Art. mit \(b^0,a^0,\lambda^0\) bezeichnet war, nemlich die bis auf die dritte Ordnung (ausschliesslich) genauen Werthe von \(b,t,l\), und zwischen denen die Gleichung
\[
rr+\tau\tau=\beta\beta+\lambda\lambda
\]
Statt findet. Die Ordnungen werden hier immer so verstanden, dass \(r\) wie eine Grösse erster Ordnung betrachtet wird.

\indent Zur Abkürzung wird noch geschrieben
\[
\begin{aligned}
\frac{2\operatorname{tang}\frac12r}{r}&=m\\
\frac{2\sin\frac12r}{r}&=n
\end{aligned}
\]
Zu der beabsichtigten Entwicklung gelangen wir am leichtesten durch Benutzung der Umwandlung der Formel
\[
x=\sin y
\]
